\section{Adaptive faithful Jones computation with a charged fallback}
\label{sec:adaptive-jones}

\subsection{A portfolio with an asymptotic guarantee}
The earlier measurements favor faithful Potts on small diagrams and tree
medials, and binary spin on larger grid diagrams. Neither implementation
dominates. The user's introsort suggestion and the certified-primitives
report's bounded-Potts-prelude proposal suggest a one-way adaptive policy:
try the practically favorable method, but bound its work before invoking
the method with the stronger general complexity guarantee.

The optional \code{faithful-adaptive} backend implements this policy for
full Jones computation and exact polynomial-identity queries. For an input
with $n$ crossings, it gives faithful Potts a default trial of $128n$
transitions. The constant is a scheduling choice, not a topological theorem.
The trial includes any internal Potts ordering restart. If Potts completes,
its faithful answer is final for this invariant query. If it hits its trial
allowance or frontier cap, the wrapper abandons the table and starts spin.
A zero trial selects spin directly. The crossing-free unknot is handled by
the existing exact empty-diagram case even with zero local caps.

This uses measured progress: the transition counter and actual state
population determine whether the prelude completes. It does not choose a
backend solely from the input crossing count or from a predicted width.
Unlike introsort, the two representations do not currently share partially
contracted coefficients. Switching restarts the invariant calculation, and
all completed work in the discarded representation is explicitly charged.

\subsection{One resource ledger across both implementations}
Let $C$ be the caller's transition allowance, possibly infinite, and $T=128n$.
The Potts attempt receives $\min(T,C)$. If it fails after $p$ transitions,
spin receives exactly $C-p$. Its completed work $s$ produces total $p+s$;
its local exhaustion reports that same cumulative count. No fresh full
allowance is allocated to the replacement. If $C-p=0$, the wrapper returns
an inconclusive limit before starting spin. An explicitly configured trial
has its own cost and does not inherit the default linear-prelude guarantee.

Both stages use the same state cap. The reported peak state count is the
maximum across attempts, including a rejected insertion that first exceeds
a state cap. It is not a byte-memory limit, nor the sum of simultaneous
populations. The failed Potts exception handler is exited before allocating
the spin table, releasing the abandoned traceback and its evaluator locals.
Only progress counters and completed order information are retained.

Potts may already have constructed a checked separator order while adapting
its own scan. The wrapper passes that order to spin without reconstructing
its certificate. Otherwise spin constructs its usual certified order.
Thus a failed prelude does not require a second certificate construction.
The retained certificate's order still has the square-root frontier bound.
Caller cancellation exceptions propagate. A local limit can retain a completed
order certificate and policy counters, but never an unfinished polynomial,
identity assertion or scalar value.

The public result preserves the winning backend's scalar format and names
its selected backend. The obstruction wrapper chooses the corresponding
hexadecimal serialization, including large quadratic coordinates and color
counts for Potts or bracket scalars for spin. CLI output includes the policy
record. Recognition preserves that record even when the Jones result is
inconclusive and the independent Khovanov calculation continues.

\subsection{Why the prelude preserves the stronger complexity bound}
With external caps and deadlines disabled, the default Potts attempt still
has only $128n$ transitions. Each transition changes the represented table
by at most one key. Thus it retains polynomially many keys; keys have at most
$O(n)$ entries, and faithful quadratic coefficients have $O(n^2)$ bits.
Even pessimistic dictionary lookup costs remain polynomial in this bounded
attempt. Input validation, ordering, at most one separator preparation,
normalization and optional coefficient recovery likewise have polynomial
bit cost. Internal reordering does not multiply the work allowance: both
attempts contribute to the same Potts transition counter.

If the prelude completes, the whole query has polynomial cost. Otherwise
the remaining spin computation has the certified $2^{O(\sqrt n)}$ bound,
including exact arithmetic and full coefficient recovery. Hence
\[
 \operatorname{poly}(n)+2^{O(\sqrt n)}=2^{O(\sqrt n)}.
\]
This guarantee is independent of whether the prelude proves useful on the
input. It applies to the default bounded trial; an arbitrarily large
user-supplied trial contributes its own additional work. Fixed local caps
can make the query inconclusive and do not turn a partial result into a proof.

Both backends compute the same faithful invariant. They cannot disagree
about a completed polynomial-one result. However, that result still does not
certify an unknot. The general subexponential bound established here is for
Jones computation, not for complete unknot recognition. The policy offers
no theorem that spin is within a constant factor of the faster backend on
every input, and the chosen transition constant is not a runtime optimum.

\subsection{Validation and measurements}
The focused tests compare both winners and direct-spin mode against an
independent whole-cube polynomial oracle on 36 diagram/mirror inputs, giving
108 full-polynomial comparisons. They verify that easy Potts completions
avoid spin and separator preparation; that a measured-work switch reuses its
certificate; that a state-limited Potts query can finish in spin; and that
exact shared-budget boundaries give completion or inconclusive exhaustion.
A shuffled grid forces a Potts ordering restart before the cross-backend
switch and verifies that both discarded stages are charged. Zero budgets,
empty input, malformed options, cancellation during either decoder, large
JSON evidence, CLI dispatch and Khovanov fallback after Jones identity are
also covered.

The integrated suite passed 725 tests in 165.183 seconds. A final refinement
counts a switch only when spin actually starts, and extends the nested
ordering-restart assertion. After that refinement, all seven focused test
methods passed again in 1.637 seconds. The broad suite predates that small
counter clarification; the final focused gate covers it explicitly.

The maintained \path{benchmark_adaptive_jones.py} compares historical Potts
at commit \code{153bc2362}, an identical A/A control, current valuation spin,
and the new adaptive wrapper. Fourteen inputs use seed 26100867 to shuffle
five measured rounds after one excluded warmup. Each query includes fresh
PD validation, ordering, tensor preparation, contraction, normalization and
full coefficient recovery. The common allowances are three cooperative
seconds, 4096 states and 200000 transitions. Source hashes are frozen across
the run. All 56 warmup attempts are excluded from timings. Of 280 measured
attempts, 270 complete; the ten incomplete attempts are the two Potts arms
on the 196-crossing grid. They are never treated as completed-query times.
Small diagrams use an independent whole-cube oracle, grids have verified
descending certificates, and the other cases compare completed polynomials
across arms. These are full invariant queries, not full recognition timings.

\begin{center}\small
\begin{tabular}{@{}lrrrr@{}}
\toprule
Input & Potts & A/A & Spin & Adaptive \\
\midrule
Trefoil & 0.185 & 0.194 & 0.589 & 0.193 \\
Conway & 0.646 & 0.656 & 2.555 & 0.671 \\
KT & 0.695 & 0.714 & 2.675 & 0.721 \\
Hard 8 & 0.360 & 0.378 & 1.382 & 0.393 \\
Braid5 36 & 3.127 & 3.141 & 13.151 & 3.149 \\
Weaving 10 & 0.631 & 0.629 & 2.508 & 0.644 \\
Torus 101 & 37.607 & 37.543 & 41.431 & 37.330 \\
Tree 6 & 10.586 & 10.680 & 43.437 & 10.660 \\
Tree 7 & 73.716 & 122.667 & 116.558 & 84.612 \\
Grid 8 & 6.740 & 6.799 & 24.626 & 6.726 \\
Grid 10 & 38.881 & 38.092 & 67.187 & 38.037 \\
Grid 12 & 367.153 & 362.538 & 176.774 & 279.949 \\
Grid 14 & limit & limit & 480.226 & 507.184 \\
Grid 8 shuffled & 28.519 & 28.516 & 24.128 & 28.537 \\
\bottomrule
\end{tabular}
\end{center}
Median milliseconds for fresh full Jones queries, including ordering. Limits are censored, not completed-query times.


On the 144-crossing grid, adaptive computation uses 18432 discarded Potts
transitions and 32098 spin transitions, totaling 50530. It reuses the checked
order. The adaptive median is 279.949 milliseconds, compared with 367.153
for Potts and 176.774 for direct spin; the median paired Potts/adaptive ratio
is 1.330. Thus the policy improves on Potts here, but the unsuccessful prelude
still has a measurable cost relative to direct spin.

On the 196-crossing grid, Potts exceeds its state cap after 6405 transitions.
Adaptive computation then spends 89122 spin transitions and completes in
507.184 milliseconds, compared with 480.226 for direct spin. The ten capped
Potts attempts establish a completion difference under the same resource
limits, not a finite completed-query speedup. On the 144-crossing case,
the focused test confirms completion at the exact combined work allowance;
one transition less gives an inconclusive result with the charged cumulative count.

Most small cases retain Potts's timing advantage over spin, with some wrapper
overhead. For example, the trefoil medians are 0.185 milliseconds for Potts,
0.589 for spin and 0.193 for adaptive computation. The 126-crossing tree has
medians 10.586, 43.437 and 10.660 respectively. The 254-crossing tree has
substantial timing variation, including different medians for the identical
Potts arms, so it does not support a precise small percentage comparison.
The shuffled 64-crossing grid finishes within the Potts prelude after an
internal ordering restart; adaptive takes 28.537 milliseconds versus 24.128
for direct spin. These counterexamples rule out interpreting the policy as
an oracle for the fastest representation. It remains optional, with the
existing overall backend default unchanged.

Raw measurements are in \path{results/adaptive_jones_20261008.json}; test
logs and PDF validation are retained under \path{synthesis/data/}. The
research refresh finds unchanged hashes for all 21 incoming archives and
unchanged \code{openai/math} head \code{fd4aeeb2ee}. This checkpoint promotes
the bounded-prelude proposal; the other component reviews in
Section~\ref{sec:incoming-oct8} remain separate work.
The subsequent intake commit \code{b71e64c3b} places those six deliveries in
\path{reports/29/} through \path{reports/34/}. Comparison against the original
archives finds only omitted checksum manifests and CRLF normalization in two
CSV files; all other placed bytes match. This placement adds no new research
claim to the current review queue.
